% JETSPIN Package documentation

\section{The JETSPIN Package\label{sec:The-JETSPIN-Package}}

In the recent years, electrospun nanofibers have gained a considerable
industrial interest due to many possible applications, such as tissue
engineering, air and water filtration, drug delivery and regenerative
medicine. In particular, the high surface-area ratio of the fibers
offers an intriguing prospect for technological applications. As consequence,
several studies were focused on the characterization and production
of uni-dimensionally elongated nanostructures. A number of reviews
\cite{li2004electrospinning,greiner2007electrospinning,carroll2008nanofibers,huang2003review,persano2013industrial}
and books \cite{yarin2014fundamentals,wendorff2012electrospinning,pisignanoelectrospinning}
concerning electrospinning have been published in the last decade.

Typically, electrospun nanofibers are produced at laboratory scale
by the uniaxial stretching of a jet, which is ejected at the nozzle
from a charged polymer solution. The initial elongation of a jet can
be produced by applying an externally electrostatic field between
the spinneret and a conductive collector. Electrospinning involves
mainly two sequential stages in the uniaxial elongation of the extruded
polymer jet: an initial quasi steady stage, in which the electric
field stretches the jet in a straight path away from the nozzle of
the ejecting apparatus, and a second stage characterized by a bending
instability induced from small perturbations, which misalign the jet
from its axis of elongation \cite{zeng2006numerical}. These small
disturbances may originate from mechanical vibrations at the nozzle
or from hydrodynamic-aerodynamic perturbations within the experimental
apparatus. Such a misalignment provides an electrostatic-driven bending
instability before the jet reaches the conductive collector, where
the fibers are finally deposited. As a consequence, the jet path length
between the nozzle and the collector increases, and the stream cross-section
undergoes a further decrease. The ultimate goal of electrospinning
process is to minimize the radius of the collected fibers. By a simple
argument of mass conservation, this is tantamount to maximizing the
jet length by the time it reaches the collecting plane. By the same
argument, it is therefore of interest to minimize the length of the
initial stable jet region. Consequently, the bending instability is
a desirable effect, as it produces a higher surface-area-to-volume
ratio of the jet, which is transferred to the resulting nanofibers
\cite{feng2003stretching}.

Computational models provide a useful tool to elucidate the physical
phenomenon and provide information which might be used for the design
of electrospinning experiments. Numerical simulations can be used
to improve the capability of predicting the key processing parameters
and exert a better control on the resulting nanofiber structure. In
recent years, with a renewed interest in nanotechnology, electrospinning
studies attracted the attention of many resarchers both from modeling,
computational and experimental point of view \cite{reneker2000bending,yarin2001taylor,fridrikh2003controlling,theron2004experimental,lu2006computer,zeng2006numerical}.
Although some author uses dissipative particle dynamics (DPD) mesoscale
simulation method into electrospinning study \cite{wang2013simulation},
most models usually treat the jet filament as a series of discrete
elements ({\em beads}) obeying the equations of continuum mechanics
\cite{reneker2000bending,yarin2001taylor}. Each bead is subject to
different types of interactions, such as long-range Coulomb repulsion,
viscoelastic drag, the external electric field. The main aim of such
models is to explain the complexity of the resulting dynamics and
provide the set of parameters driving the optimal process. The effect
of fast-oscillating loads on the bending instability have been explored
in an extensive modeling and computational study \cite{coluzza2014ultrathin}.

JETSPIN delivers a FORTRAN code especially designed to simulate the
electrospinning process in a variety of different conditions and experimental
settings. This comprehensive platform is devised such that different
cases and input variables can be described and simulated. The framework
is developed to exploit several computational architectures, both
serial and parallel.

JETSPIN, as open-source software, can be used to carry out a systematic
sensitivity analysis over a broad range of parameter values. The results
of simulations provide valuable insight on the physics of the process
and can be used to assess experimental procedures for an optimal design
of the equipment and to control processing strategies of technological
advanced nanofibers.

\textbf{The JETSPIN code is freely downloadable at the web site: }

\bigskip{}

\begin{center}
\href{http://www.nanojets.eu}{http://www.nanojets.eu}
\par\end{center}
